\chapter{Probability}

\section{Sample spaces and events}

\begin{definition}
An \emph{experiment} is a process with uncertain outcome. The
\emph{sample space} $\Omega$ is the set of all possible outcomes. An
\emph{event} is a subset $A\subseteq\Omega$.
\end{definition}

\begin{example}
A susceptible individual has one contact with an infectious person.
\[
\Omega=\{\text{infected},\ \text{not infected}\}.
\]
$A=\{\text{infected}\}$ is an event.
\end{example}

\section{Set operations on events}

\begin{itemize}
  \item \textbf{Union} $A\cup B$: ``$A$ or $B$ occurs.''
  \item \textbf{Intersection} $A\cap B$: ``$A$ and $B$ occur.''
  \item \textbf{Complement} $A^c$: ``$A$ does not occur.''
  \item \textbf{Disjoint}: $A\cap B=\emptyset$ ($A$ and $B$ cannot occur together).
\end{itemize}

\begin{example}
$A=\{$infected with strain 1$\}$, $B=\{$infected with strain 2$\}$.
\begin{itemize}[nosep]
  \item $A\cup B$: infected with at least one strain.
  \item $A\cap B$: co-infected with both strains.
  \item $A^c$: not infected with strain 1.
\end{itemize}
\end{example}

\section{The probability function}

\begin{definition}[Axioms of probability]
$\Pr$ assigns to each event $A\subseteq\Omega$ a number $\Pr(A)\in[0,1]$ such that
\begin{itemize}[nosep]
  \item $\Pr(\Omega)=1$;
  \item for disjoint events $A,B$, $\ \Pr(A\cup B)=\Pr(A)+\Pr(B)$.
\end{itemize}
\end{definition}

Interpretation: $\Pr(A)$ is the long-run fraction of times $A$ occurs (e.g. $\Pr(\text{infection})=0.3$ means 30\% of exposed contacts result in infection).

\section{Probability of the complement}
A very useful identity is $\Pr(A^c) = 1-\Pr(A)$.

%% \begin{proposition}[Addition law]
%% $\Pr(A\cup B) = \Pr(A)+\Pr(B)-\Pr(A\cap B)$.
%% \end{proposition}

\section{Conditional probabilities}
\[
\Pr(A|B) := \frac{\Pr(A \cap B)}{\Pr(B)} \qquad \text{whenever $\Pr(B) > 0$.}
\]
\section{Independence}

\begin{definition}
Events $A$ and $B$ are \emph{independent} if
\[
\Pr(A \cap B) = \Pr(A)\,\Pr(B).
\]
\end{definition}

If moreover $\Pr(B)>0$, this is equivalent (by the definition of
conditional probability above) to $\Pr(A|B) = \Pr(A)$.

Interpretation: knowing that some of the events occurred gives no
information about whether the others occurred.

\begin{example}
A susceptible individual has $n$ contacts. Each contact independently
transmits infection with probability $c$. Let $A_i=\{\text{contact $i$
  transmits infection}\}$, so $\Pr(A_i)=q$. Avoiding infection on
all $n$ contacts is the event $A_1^c\cap\cdots\cap A_n^c$, so by
independence
\[
\Pr(\text{avoid infection}) = \prod_{i=1}^n \Pr(A_i^c) = (1-q)^n.
\]
\end{example}

\section{Random variables}

\begin{definition}
A \emph{random variable} $X$ is a function $X:\Omega\to\mathbb{R}$ assigning a number to each outcome.
\end{definition}

\begin{example}
$T=$ time until recovery.
\end{example}

\begin{definition}
The \emph{density} (pdf) of $X$ is a function $f\ge 0$ with $\displaystyle\int_{-\infty}^\infty f(x)\,dx=1$ such that
\[
\Pr(a\le X\le b)=\int_a^b f(x)\,dx.
\]
The \emph{expectation} is $E[X]=\int x f(x)\,dx$.
\end{definition}

\section{The exponential random variable}

Used to model waiting times: time to recovery, time until next contact, incubation period.

\begin{definition}
$T$ has an \emph{exponential distribution with rate $\lambda>0$}, written $T\sim\text{Exp}(\lambda)$, if
\[
f(t)=\lambda e^{-\lambda t}, \qquad t\ge 0.
\]
\end{definition}

The \emph{survival function} is the probability of not yet having recovered by time $t$:
\[
S(t) = \Pr(T>t) = \int_t^\infty \lambda e^{-\lambda s}\,ds = e^{-\lambda t}.
\]
Mean: $E[T]=1/\lambda$.

Interpretation: if $T$ is the infectious period and $T\sim\text{Exp}(\gamma)$, then $\gamma$ is the recovery \emph{rate} and $1/\gamma$ is the mean infectious period.

\section{The memoryless property}
\begin{proposition}
  If $T\sim\text{Exp}(\lambda)$, then $T$ is memoryless, i.e.~
  \[
  \Pr(T> t+s \mid T> t) = \Pr(T> s) \qquad \text{for all } s,t\ge 0.
  \]
\end{proposition}
Interpretation: given that an individual is still infectious at time
$t$, the probability they remain infectious for $s$ more days does not
depend on $t$. Time already spent infected gives no information about
time remaining --- the recovery rate is constant, with no ``aging'' or
``wearing out.''

\begin{proof}
\begin{align*}
  \Pr(T> t+s\mid T> t) &= \frac{\Pr(T> t+s \cap T > t)}{\Pr(T > t)}\\
  &= \frac{\Pr(T> t+s)}{\Pr(T> t)}\\
  &= \frac{e^{-\lambda(t+s)}}{e^{-\lambda t}}\\
  &= e^{-\lambda s} \\
  &= \Pr(T> s).
\end{align*}
\end{proof}

\section{Minimum of independent exponentials}

\begin{proposition}
  If $X\sim\text{Exp}(\lambda)$ and $Y\sim\text{Exp}(\eta)$ are
  independent, define $Z=\min\{X,Y\}$.
\[
Z\sim\text{Exp}(\lambda+\eta).
\]
\end{proposition}

Interpretation: if an infected individual leaves the infected class
either by recovering (rate $\gamma$) or by dying of unrelated causes
(rate $\mu$), independently, then the time spent in the infected class
is $\min(T_{\text{recovery}}, T_{\text{death}})\sim\text{Exp}(\gamma+\mu)$
--- whichever happens first --- with mean $1/(\gamma+\mu)$.

\begin{proof}
By independence
\begin{align*}
  \Pr(Z>t) &= \Pr(X>t)\,\Pr(Y>t) \\
  &= e^{-\lambda t}e^{-\eta t}\\
  &= e^{-(\lambda+\eta)t} \\
\end{align*}
which is indeed the survival function of $\text{Exp}(\lambda+\eta)$.
\end{proof}
